% Options for packages loaded elsewhere
\PassOptionsToPackage{unicode}{hyperref}
\PassOptionsToPackage{hyphens}{url}
\PassOptionsToPackage{table}{xcolor}
%
\documentclass[
  letterpaper,
]{article}
\usepackage{amsmath,amssymb}
\usepackage{iftex}
\usepackage[T1]{fontenc}
\usepackage[utf8]{inputenc}
\usepackage{textcomp}
\usepackage{lmodern}
\ifPDFTeX\else
  % xetex/luatex font selection
\fi
% Use upquote if available, for straight quotes in verbatim environments
\IfFileExists{upquote.sty}{\usepackage{upquote}}{}
\IfFileExists{microtype.sty}{% use microtype if available
  \usepackage[]{microtype}
  \UseMicrotypeSet[protrusion]{basicmath} % disable protrusion for tt fonts
}{}
\makeatletter
\@ifundefined{KOMAClassName}{% if non-KOMA class
  \IfFileExists{parskip.sty}{%
    \usepackage{parskip}
  }{% else
    \setlength{\parindent}{0pt}
    \setlength{\parskip}{6pt plus 2pt minus 1pt}}
}{% if KOMA class
  \KOMAoptions{parskip=half}}
\makeatother
\usepackage{xcolor}
\usepackage{color}
\usepackage{fancyvrb}
\newcommand{\VerbBar}{|}
\newcommand{\VERB}{\Verb[commandchars=\\\{\}]}
\DefineVerbatimEnvironment{Highlighting}{Verbatim}{commandchars=\\\{\}}
% Add ',fontsize=\small' for more characters per line
\newenvironment{Shaded}{}{}
\newcommand{\AlertTok}[1]{\textcolor[rgb]{1.00,0.00,0.00}{\textbf{#1}}}
\newcommand{\AnnotationTok}[1]{\textcolor[rgb]{0.38,0.63,0.69}{\textbf{\textit{#1}}}}
\newcommand{\AttributeTok}[1]{\textcolor[rgb]{0.49,0.56,0.16}{#1}}
\newcommand{\BaseNTok}[1]{\textcolor[rgb]{0.25,0.63,0.44}{#1}}
\newcommand{\BuiltInTok}[1]{\textcolor[rgb]{0.00,0.50,0.00}{#1}}
\newcommand{\CharTok}[1]{\textcolor[rgb]{0.25,0.44,0.63}{#1}}
\newcommand{\CommentTok}[1]{\textcolor[rgb]{0.38,0.63,0.69}{\textit{#1}}}
\newcommand{\CommentVarTok}[1]{\textcolor[rgb]{0.38,0.63,0.69}{\textbf{\textit{#1}}}}
\newcommand{\ConstantTok}[1]{\textcolor[rgb]{0.53,0.00,0.00}{#1}}
\newcommand{\ControlFlowTok}[1]{\textcolor[rgb]{0.00,0.44,0.13}{\textbf{#1}}}
\newcommand{\DataTypeTok}[1]{\textcolor[rgb]{0.56,0.13,0.00}{#1}}
\newcommand{\DecValTok}[1]{\textcolor[rgb]{0.25,0.63,0.44}{#1}}
\newcommand{\DocumentationTok}[1]{\textcolor[rgb]{0.73,0.13,0.13}{\textit{#1}}}
\newcommand{\ErrorTok}[1]{\textcolor[rgb]{1.00,0.00,0.00}{\textbf{#1}}}
\newcommand{\ExtensionTok}[1]{#1}
\newcommand{\FloatTok}[1]{\textcolor[rgb]{0.25,0.63,0.44}{#1}}
\newcommand{\FunctionTok}[1]{\textcolor[rgb]{0.02,0.16,0.49}{#1}}
\newcommand{\ImportTok}[1]{\textcolor[rgb]{0.00,0.50,0.00}{\textbf{#1}}}
\newcommand{\InformationTok}[1]{\textcolor[rgb]{0.38,0.63,0.69}{\textbf{\textit{#1}}}}
\newcommand{\KeywordTok}[1]{\textcolor[rgb]{0.00,0.44,0.13}{\textbf{#1}}}
\newcommand{\NormalTok}[1]{#1}
\newcommand{\OperatorTok}[1]{\textcolor[rgb]{0.40,0.40,0.40}{#1}}
\newcommand{\OtherTok}[1]{\textcolor[rgb]{0.00,0.44,0.13}{#1}}
\newcommand{\PreprocessorTok}[1]{\textcolor[rgb]{0.74,0.48,0.00}{#1}}
\newcommand{\RegionMarkerTok}[1]{#1}
\newcommand{\SpecialCharTok}[1]{\textcolor[rgb]{0.25,0.44,0.63}{#1}}
\newcommand{\SpecialStringTok}[1]{\textcolor[rgb]{0.73,0.40,0.53}{#1}}
\newcommand{\StringTok}[1]{\textcolor[rgb]{0.25,0.44,0.63}{#1}}
\newcommand{\VariableTok}[1]{\textcolor[rgb]{0.10,0.09,0.49}{#1}}
\newcommand{\VerbatimStringTok}[1]{\textcolor[rgb]{0.25,0.44,0.63}{#1}}
\newcommand{\WarningTok}[1]{\textcolor[rgb]{0.38,0.63,0.69}{\textbf{\textit{#1}}}}
\usepackage{longtable,booktabs,array}
\usepackage{calc} % for calculating minipage widths
% Correct order of tables after \paragraph or \subparagraph
\usepackage{etoolbox}
\makeatletter
\patchcmd\longtable{\par}{\if@noskipsec\mbox{}\fi\par}{}{}
\makeatother
% Allow footnotes in longtable head/foot
\IfFileExists{footnotehyper.sty}{\usepackage{footnotehyper}}{\usepackage{footnote}}
\makesavenoteenv{longtable}
\setlength{\emergencystretch}{3em} % prevent overfull lines
\providecommand{\tightlist}{%
  \setlength{\itemsep}{0pt}\setlength{\parskip}{0pt}}
\setcounter{secnumdepth}{5}
% Self-contained presentation rules embedded in every generated LaTeX file.
\usepackage{iftex}
\ifPDFTeX\else\PackageError{SSE-ASAS}{Compile with pdfLaTeX only}{}\fi
\usepackage[T1]{fontenc}
\usepackage[utf8]{inputenc}
\usepackage{lmodern}
\usepackage{microtype}
\usepackage{geometry}
\geometry{margin=0.93in,headheight=16pt}
\usepackage{xcolor} % [table] passed at top of file (option-clash fix)
\definecolor{ink}{HTML}{15283B}
\definecolor{brand}{HTML}{075B75}
\definecolor{brandlight}{HTML}{EAF5F8}
\definecolor{warm}{HTML}{A15C1D}
\definecolor{muted}{HTML}{536577}
\usepackage{fancyhdr}
\usepackage{titlesec}
\usepackage{enumitem}
\usepackage{booktabs}
\usepackage{longtable}
\usepackage{array}
\usepackage{tabularx}
\usepackage{tcolorbox}
\tcbuselibrary{breakable,skins}
\usepackage{listings}
\usepackage{fancyvrb}
\usepackage{hyperref}
\hypersetup{colorlinks=true,linkcolor=brand,urlcolor=brand,citecolor=brand}
\urlstyle{same}
\usepackage{lastpage}
\setlength{\emergencystretch}{3em}
\setlength{\parskip}{0.45em}
\setlength{\parindent}{0pt}
\setlist[itemize]{leftmargin=*,topsep=3pt,itemsep=2pt}
\setlist[enumerate]{leftmargin=*,topsep=3pt,itemsep=2pt}
\titleformat{\section}{\Large\bfseries\color{ink}}{\thesection}{0.75em}{}[\vspace{-2pt}\color{brand}\titlerule]
\titleformat{\subsection}{\large\bfseries\color{brand}}{\thesubsection}{0.7em}{}
\titleformat{\subsubsection}{\normalsize\bfseries\color{ink}}{\thesubsubsection}{0.7em}{}
\pagestyle{fancy}
\fancyhf{}
\fancyhead[L]{\small\sffamily\textcolor{brand}{SOFTWARE SECURITY ENGINEERING}}
\fancyhead[R]{\small\sffamily\textcolor{muted}{TECHNICAL SERIES}}
\fancyfoot[L]{\scriptsize\sffamily Version 1.0 \textbullet\ 14 October 2026}
\fancyfoot[R]{\scriptsize\sffamily Page \thepage\ of \pageref{LastPage}}
\renewcommand{\headrulewidth}{0.5pt}
\renewcommand{\footrulewidth}{0.2pt}
\setcounter{tocdepth}{2}
\setcounter{secnumdepth}{3}
\lstdefinestyle{sse}{basicstyle=\ttfamily\footnotesize,frame=single,framerule=0.4pt,rulecolor=\color{brand!40},backgroundcolor=\color{brandlight!45},numbers=left,numberstyle=\scriptsize\color{muted},numbersep=6pt,breaklines=true,breakatwhitespace=false,tabsize=2,showstringspaces=false,columns=fullflexible,keepspaces=true,captionpos=b,aboveskip=9pt,belowskip=9pt}
\lstset{style=sse}
\newtcolorbox{sourceconcept}[1][]{enhanced,breakable,colback=brandlight,colframe=brand!65!black,title={ORIGINAL CONCEPT},fonttitle=\bfseries\sffamily,#1}
\newtcolorbox{modernextension}[1][]{enhanced,breakable,colback=white,colframe=warm!85!black,title={MODERN EXTENSION},fonttitle=\bfseries\sffamily,#1}
\newtcolorbox{fieldwork}[1][]{enhanced,breakable,colback=white,colframe=ink!55,title={PRACTICAL APPLICATION},fonttitle=\bfseries\sffamily,#1}
\AtBeginDocument{\color{ink}}
\makeatletter
\renewcommand{\maketitle}{%
 \begin{titlepage}\thispagestyle{empty}
 \vspace*{1.5cm}{\color{brand}\rule{\textwidth}{3pt}}\par\vspace{1cm}
 {\sffamily\bfseries\color{muted}\large SOFTWARE SECURITY ENGINEERING\\ APPLICATION SECURITY ASSESSMENT SERIES}\par\vspace{1.05cm}
 {\sffamily\bfseries\color{ink}\fontsize{27}{34}\selectfont\@title\par}\vspace{0.8cm}
 {\sffamily\large\color{brand}\@author\par}\vfill
 \begin{tcolorbox}[colback=brandlight,colframe=brand!55,boxrule=0.5pt]
 \small\textbf{Editorial policy}\quad Original teaching synthesis with book-grounded concepts distinguished from contemporary practices. Modern platform behavior is version-sensitive; validate against linked vendor documentation before production use.
 \end{tcolorbox}
 {\color{brand}\rule{\textwidth}{2pt}}\par\vspace{0.3cm}
 {\sffamily\small Version 1.0 \quad|\quad October 2026 \quad|\quad pdfLaTeX}\end{titlepage}\setcounter{page}{1}}
\makeatother
\ifLuaTeX
  \usepackage{selnolig}  % disable illegal ligatures
\fi
\usepackage{bookmark}
\IfFileExists{xurl.sty}{\usepackage{xurl}}{} % add URL line breaks if available
\urlstyle{same}
\hypersetup{
  pdftitle={Prerequisite-based 16-week Curriculum},
  pdfauthor={Independent Technical Study Edition},
  hidelinks,
  pdfcreator={LaTeX via pandoc}}

\title{Prerequisite-based 16-week Curriculum}
\author{Independent Technical Study Edition}
\date{}

\begin{document}
\maketitle

{
\setcounter{tocdepth}{3}
\tableofcontents
}
\section{Program charter and
outcomes}\label{program-charter-and-outcomes}

A \textbf{16-week, 7--10 hour-per-week} curriculum develops security
foundations, source auditing, SAST engineering, secure coding, and
end-to-end assessment competence. The central teaching strategy is to
preserve the theory and auditing methodologies of three foundational
texts and add a visibly separate modern toolchain. An experienced
developer may accelerate exercises after demonstrating the exit
criteria; the dependencies still apply.

\textbf{Texts:} {[}B{]} Matt Bishop, \emph{Computer Security: Art and
Science}, second edition (2019); {[}D{]} Mark Dowd, John McDonald,
Justin Schuh, \emph{The Art of Software Security Assessment} (2007);
{[}C{]} Brian Chess and Jacob West, \emph{Secure Programming with Static
Analysis} (2007). These are recommended source chapters rather than
reproductions. Their historical platforms and frameworks remain
instructive but should not be mistaken for current production
recommendations.

\textbf{Modern references} {[}M{]} are authoritative online sources:
MITRE CWE, OWASP ASVS/Top 10/SAMM, NIST SSDF, and GitHub Code
Security/CodeQL documentation. These are course extensions, not
assertions that the original books describe today's products.

\section{Prerequisite graph}\label{prerequisite-graph}

\begin{Shaded}
\begin{Highlighting}[]
\NormalTok{Programming fundamentals + Git/CLI}
\NormalTok{              |}
\NormalTok{              v}
\NormalTok{F1: Security properties, policy and assurance (W1{-}2)}
\NormalTok{              |}
\NormalTok{              v}
\NormalTok{F2: Access control, information flow, trust (W3{-}4)}
\NormalTok{              |}
\NormalTok{              v}
\NormalTok{V1: Threat modeling and vulnerability classes (W5{-}6)}
\NormalTok{              |}
\NormalTok{        +{-}{-}{-}{-}{-}+{-}{-}{-}{-}{-}{-}{-}+}
\NormalTok{        |             |}
\NormalTok{        v             v}
\NormalTok{M1: Manual audit   S1: Analysis foundations}
\NormalTok{    (W7{-}8)             (W9{-}10)}
\NormalTok{        |             |}
\NormalTok{        +{-}{-}{-}{-}{-}{-}v{-}{-}{-}{-}{-}{-}+}
\NormalTok{               |}
\NormalTok{    P1: Secure coding and testing (W11{-}12)}
\NormalTok{               |}
\NormalTok{               v}
\NormalTok{    A1: Assessment and GHAS integration (W13{-}14)}
\NormalTok{               |}
\NormalTok{               v}
\NormalTok{       C1: Capstone + evidence review (W15{-}16)}
\end{Highlighting}
\end{Shaded}

The branch after V1 reflects independent \emph{manual} and \emph{static}
investigative techniques; P1 requires both. Actual calendar sequence can
interleave two branches after their common prerequisites, but neither
substitutes for the other.

\section{Module contracts and source
mapping}\label{module-contracts-and-source-mapping}

\begin{longtable}[]{@{}
  >{\raggedright\arraybackslash}p{(\columnwidth - 6\tabcolsep) * \real{0.2500}}
  >{\raggedright\arraybackslash}p{(\columnwidth - 6\tabcolsep) * \real{0.2500}}
  >{\raggedright\arraybackslash}p{(\columnwidth - 6\tabcolsep) * \real{0.2500}}
  >{\raggedright\arraybackslash}p{(\columnwidth - 6\tabcolsep) * \real{0.2500}}@{}}
\toprule\noalign{}
\begin{minipage}[b]{\linewidth}\raggedright
ID / weeks
\end{minipage} & \begin{minipage}[b]{\linewidth}\raggedright
Competency and dependency
\end{minipage} & \begin{minipage}[b]{\linewidth}\raggedright
Primary documents
\end{minipage} & \begin{minipage}[b]{\linewidth}\raggedright
Source sections
\end{minipage} \\
\midrule\noalign{}
\endhead
\bottomrule\noalign{}
\endlastfoot
F1 (1--2) & CIA, policy, threat assumptions, assurance; general
programming & Foundations; Labs L01 & {[}B{]} 1--4, 19 \\
F2 (3--4) & Security models, reference monitor, information flow; F1 &
Foundations; L01--L03 & {[}B{]} 5--9, 14--18, 20--22 \\
V1 (5--6) & Taxonomy, trust boundaries, root causes; F2 & Taxonomy; L02,
L04 & {[}D{]} 1--3, 5--8; {[}C{]} 1, 5--8 \\
M1 (7--8) & Plan--Work--Reflect, audit trace; V1 & Manual Auditing;
L05--L07 & {[}D{]} 2--4, 9--13, 17--18 \\
S1 (9--10) & CFG, SSA, dataflow, taint, rules; V1 & Static Analysis;
L08--L10 & {[}C{]} 2--4, 13--14 \\
P1 (11--12) & Multi-language safe APIs, test evidence; M1+S1 & Secure
Coding; L11--L12 & {[}C{]} 5--12; {[}D{]} 5--13 \\
A1 (13--14) & Assessment, GitHub workflow, dispositions; P1 &
Procedures; L10, L13 & {[}D{]} 2--4; {[}B{]} 19--25; {[}C{]} 3--4 \\
C1 (15--16) & Independent capstone report; A1 & Labs L14; all guides &
Integrated synthesis {[}B{]}{[}C{]}{[}D{]}{[}M{]} \\
\end{longtable}

\section{Week-by-week teaching plan}\label{week-by-week-teaching-plan}

\subsection{Phase I: formal security reasoning (Weeks
1--4)}\label{phase-i-formal-security-reasoning-weeks-14}

\textbf{Week 1 -- Security vocabulary.} Read {[}B{]} Chapter 1 and
Foundations sections on assets, CIA, policy, trust. Deliver a property
inventory for a multi-tenant service and three observable failure
states. \emph{Checkpoint:} separate confidentiality, integrity,
availability, and accountability.

\textbf{Week 2 -- Protection states.} Read {[}B{]} Chapters 2--4. Build
access matrix L01 and show how a grant/revoke changes policy state.
\emph{Checkpoint:} reject permissions by default, including cross-tenant
object access.

\textbf{Week 3 -- Policy models.} Read {[}B{]} Chapters 5--9. Compare
Bell--LaPadula, Biba, Clark--Wilson, and history-aware conflicts.
Complete L03 business transaction model. \emph{Checkpoint:} distinguish
confidentiality policy from transactional integrity.

\textbf{Week 4 -- Mechanisms and assurance.} Read {[}B{]} Chapters
14--18 and 19--22 (selected sections). Complete L02 flow mapping and an
assurance argument. \emph{Checkpoint:} identify an enforcement-bypass
path and propose independent evidence.

\subsection{Phase II: vulnerabilities and manual assessment (Weeks
5--8)}\label{phase-ii-vulnerabilities-and-manual-assessment-weeks-58}

\textbf{Week 5 -- Vulnerability fundamentals.} Read {[}D{]} Chapters
1--3 and Taxonomy. Trace data, trust, and interpretation boundaries.
Compare lifecycle origin with security property. \emph{Checkpoint:}
classify five issues without confusing a CVE with a CWE.

\textbf{Week 6 -- Vulnerability mechanics.} Read {[}D{]} Chapters 5--8
and {[}C{]} Chapters 5--8; complete L04. \emph{Checkpoint:} reason
accurately about signed/unsigned conversion, SQL interpretation, and
concurrency conditions.

\textbf{Week 7 -- Audit preparation.} Read {[}D{]} Chapter 4, especially
preassessment and Plan--Work--Reflect; complete L05. \emph{Checkpoint:}
produce entry-point inventory, audit hypotheses, and scope-limited
trace.

\textbf{Week 8 -- Manual audit.} Complete L06 and L07, review {[}D{]}
Chapters 9--13, 17--18 selectively. \emph{Checkpoint:} demonstrate
object authorization fix and safe size arithmetic, with tests.

\subsection{Phase III: static analysis and defensive code (Weeks
9--12)}\label{phase-iii-static-analysis-and-defensive-code-weeks-912}

\textbf{Week 9 -- CFG and data flow.} Read {[}C{]} Chapters 2 and 4;
complete L08. \emph{Checkpoint:} compute taint facts at each CFG node
and explain fixpoint convergence.

\textbf{Week 10 -- Rules, models, and reports.} Read {[}C{]} Chapters
3--4; complete L09, start L10. \emph{Checkpoint:} show source,
propagator, sink, sanitizer and one known false-negative scenario.

\textbf{Week 11 -- Safe implementation.} Read {[}C{]} Chapters 5--9;
work Secure Coding Python/C/Java/JS examples; complete L11.
\emph{Checkpoint:} a negative regression test shows why a fix is
effective.

\textbf{Week 12 -- Data, privilege, secrets.} Read {[}C{]} Chapters
10--12 and applicable {[}D{]} platform sections; complete L12.
\emph{Checkpoint:} propose coordinated credential response,
resource/exception safety, and supporting test evidence.

\subsection{Phase IV: operational AppSec and assessment (Weeks
13--16)}\label{phase-iv-operational-appsec-and-assessment-weeks-1316}

\textbf{Week 13 -- Assessment design.} Read {[}D{]} Chapters 2--4; use
Procedures to produce an authorized charter, ASVS-informed test matrix,
and coverage declaration. \emph{Checkpoint:} scan success and
analyzed-source coverage are verified independently.

\textbf{Week 14 -- CodeQL/GitHub and triage.} Complete L10 and L13.
Validate default or advanced code scanning choice, PR workflow/branch
associations, SARIF trace, and a report that distinguishes rule
confidence from business impact. \emph{Checkpoint:} unresolved issues
have owners, evidence, and valid dispositions.

\textbf{Week 15 -- Capstone execution.} Start L14 on synthetic, local
code; submit threat model, source inventory, review plan, and first
evidence packet. \emph{Checkpoint:} peer can reproduce a finding using
its source revision and conditions.

\textbf{Week 16 -- Remediation and defense.} Close L14 with patch,
negative regression, PR scan evidence or explicit limitation, formal
residual risk, and executive summary. \emph{Checkpoint:} oral defense of
one false positive, one verified true defect, and one area of unknown
coverage.

\section{Workload and evidence map}\label{workload-and-evidence-map}

\begin{longtable}[]{@{}
  >{\raggedright\arraybackslash}p{(\columnwidth - 4\tabcolsep) * \real{0.3333}}
  >{\raggedleft\arraybackslash}p{(\columnwidth - 4\tabcolsep) * \real{0.3333}}
  >{\raggedright\arraybackslash}p{(\columnwidth - 4\tabcolsep) * \real{0.3333}}@{}}
\toprule\noalign{}
\begin{minipage}[b]{\linewidth}\raggedright
Activity
\end{minipage} & \begin{minipage}[b]{\linewidth}\raggedleft
Weekly average
\end{minipage} & \begin{minipage}[b]{\linewidth}\raggedright
Examples of tangible evidence
\end{minipage} \\
\midrule\noalign{}
\endhead
\bottomrule\noalign{}
\endlastfoot
Directed textbook reading & 2--3 h & Structured technical notes and
concept comparisons \\
Guided document study & 1--2 h & Reusable policy, audit, algorithm, or
procedure worksheet \\
Lab and implementation & 3--4 h & Code, tests, diagrams, tool outputs,
scope log \\
Reflection and review & 1 h & Finding packet and limitations ledger \\
\end{longtable}

\textbf{Milestones:} end Week 4 formal policy portfolio; end Week 8
vulnerability audit packet; end Week 10 static analysis artifact; end
Week 12 secure-coding portfolio; end Week 14 assessment plan and
scanning evidence; end Week 16 capstone acceptance.

\section{Scoring, prerequisite exemptions, and
mastery}\label{scoring-prerequisite-exemptions-and-mastery}

\begin{longtable}[]{@{}
  >{\raggedright\arraybackslash}p{(\columnwidth - 4\tabcolsep) * \real{0.3333}}
  >{\raggedleft\arraybackslash}p{(\columnwidth - 4\tabcolsep) * \real{0.3333}}
  >{\raggedright\arraybackslash}p{(\columnwidth - 4\tabcolsep) * \real{0.3333}}@{}}
\toprule\noalign{}
\begin{minipage}[b]{\linewidth}\raggedright
Competency
\end{minipage} & \begin{minipage}[b]{\linewidth}\raggedleft
Weight
\end{minipage} & \begin{minipage}[b]{\linewidth}\raggedright
Mastery evidence
\end{minipage} \\
\midrule\noalign{}
\endhead
\bottomrule\noalign{}
\endlastfoot
Formal security modeling & 15\% & Precise invariants, policy matrix,
assurance \\
Vulnerability mechanics & 15\% & Root cause and justified CWE mapping \\
Manual auditing & 20\% & Reproducible trace and design review \\
SAST algorithms and tooling & 20\% & CFG/fixpoint and accurate tool
scope \\
Secure coding and testing & 15\% & Patch plus
failing-before/passing-after test \\
Assessment governance & 15\% & Scope, reporting, residual risk, decision
audit \\
\end{longtable}

Require at least \textbf{80\% overall} and \textbf{70\% in every
competency}, plus an accepted capstone with no material safety or
evidence-integrity defects. A prior-experience exemption requires the
same assessed exit artifact; job title or certificate alone is
insufficient evidence of specific ability.

\section{Pathways and electives}\label{pathways-and-electives}

\textbf{Application security engineer:} emphasize M1/S1 and A1; add
authorization matrices, multi-tenant testing, query modeling, SARIF and
GitHub branch protection. \textbf{Native/software systems assessor:}
deepen memory semantics, concurrency, OS privilege, RPC and parsers from
{[}D{]} Chapters 5--16. \textbf{Security architecture reviewer:} deepen
{[}B{]} confidentiality/integrity/composition/assurance and {[}D{]}
design review. \textbf{DevSecOps engineer:} deepen scan coverage,
ruleset management, dependency and secret posture, policy gates, and
SSDF/SAMM maturity assessment. Electives may run after core mastery;
none replaces the required capstone.

\section{Distinguishing historical material and current
guidance}\label{distinguishing-historical-material-and-current-guidance}

\textbf{Preserve the originals:} formal protection models and assurance
{[}B{]}; a disciplined manual audit process and deep failure mechanics
{[}D{]}; static analysis representations, rule types, and reporting
lessons {[}C{]}. Historical C/Windows/UNIX/API cases remain instructive
to explain a mechanism, not to prescribe an outdated API in production.

\textbf{Modernize explicitly:} current CWE definitions; OWASP ASVS 5.0.0
and Top 10:2025; NIST SSDF 1.1; CodeQL and GitHub Code Security/Secret
Protection; SARIF workflows; contemporary language APIs and CI/CD. The
edition/version matters when interpreting specific controls and feature
behavior. Check the linked authoritative sources at implementation time.

\section{Course evaluation and
maintenance}\label{course-evaluation-and-maintenance}

At the end of each cohort, collect anonymous rubric scores and
categorize recurring errors: policy/mechanism confusion, invalid CWE
mapping, missing path conditions, scanner-output overclaim, weak
regression evidence, and unauthorized testing risk. Update
\textbf{modern-only sections} quarterly, recording verified
documentation date and changed behavior. Keep source-concept chapters
stable except for clarifications explicitly labeled editorial.
Regenerate PDFs using pdfLaTeX and check page geometry, links, tables,
and code listings.

\section{Annex A: competency-to-deliverable
traceability}\label{annex-a-competency-to-deliverable-traceability}

\subsection{Formal security knowledge (F1 and
F2)}\label{formal-security-knowledge-f1-and-f2}

The learner must differentiate policy, mechanism, and assurance in an
unfamiliar system, not merely memorize definitions. In the evaluation,
give them an approval service with worker and administrative paths and
ask them to state an integrity invariant. Strong evidence names every
actor, object and state transition, identifies a trusted enforcement
point, and provides one counterexample. Success also requires
identifying why role checks alone are insufficient for cross-tenant
access.

\subsection{Vulnerability engineering
(V1)}\label{vulnerability-engineering-v1}

Provide a snippet with two independent weaknesses, such as SQL string
composition and missing record ownership. A learner who fixes only the
SQL parser still fails the access-control requirement. They must locate
both roots, associate the relevant security properties, provide
candidate CWEs, and demonstrate a negative test for each. The exercise
guards against equating a scanner's weakness vocabulary with the
complete application policy.

\subsection{Manual investigation (M1)}\label{manual-investigation-m1}

Provide a partial call graph where a public route delegates to a service
and a worker. The learner must use a Plan--Work--Reflect loop with one
candidate-point strategy and one architecture-led strategy, explaining
findings and unknowns. Good review notes include source revision,
evidence, contradictions, and a decision about whether to continue
reviewing a sibling API.

\subsection{Static analysis (S1)}\label{static-analysis-s1}

A learner should construct an elementary forward may-taint analysis,
explain branch joins, model a context-specific sanitizer, and show how
an unmodeled framework API produces a false negative. For CodeQL
operational competence, verify a successful scan and its extraction
scope instead of accepting a green check that did no analysis. They must
be able to describe when custom models or manual review are required.

\subsection{Engineering and governance
(P1/A1/C1)}\label{engineering-and-governance-p1a1c1}

The learner patches a flaw, adds both allowed and denied behavior tests,
tracks its review in a change-controlled repository, and produces a
coherent closure record. Capstone assessment is evaluated as a product:
the artifacts should help another assessor repeat the reasoning and a
developer reproduce the fix, while allowing an authorized risk owner to
assess unresolved exposure.

\section{Annex B: deeper optional reading
tracks}\label{annex-b-deeper-optional-reading-tracks}

\begin{longtable}[]{@{}
  >{\raggedright\arraybackslash}p{(\columnwidth - 4\tabcolsep) * \real{0.3333}}
  >{\raggedright\arraybackslash}p{(\columnwidth - 4\tabcolsep) * \real{0.3333}}
  >{\raggedright\arraybackslash}p{(\columnwidth - 4\tabcolsep) * \real{0.3333}}@{}}
\toprule\noalign{}
\begin{minipage}[b]{\linewidth}\raggedright
Track
\end{minipage} & \begin{minipage}[b]{\linewidth}\raggedright
Expanded reading focus
\end{minipage} & \begin{minipage}[b]{\linewidth}\raggedright
Additional assessment challenge
\end{minipage} \\
\midrule\noalign{}
\endhead
\bottomrule\noalign{}
\endlastfoot
Formal assurance & {[}B{]} Chapters 3, 9, 19--22 & Formalize an
information-flow invariant and identify assumption gaps \\
Unix security & {[}D{]} Chapters 9--10; {[}B{]} 16--18 & Review secure
file operation and privilege transitions \\
Windows security & {[}D{]} Chapters 11--12 & Trace token, object,
impersonation, and IPC boundaries \\
Native concurrency & {[}D{]} Chapter 13; {[}C{]} Chapters 6--8 & Show a
race/ownership invariant with deterministic tests \\
Web and protocols & {[}D{]} Chapters 14, 16--18; {[}C{]} 9--10 & Review
authorization across API, parser, gateway and worker \\
Static-analysis R\&D & {[}C{]} Chapter 4 and exercises 13--14 & Extend
toy analyzer with aliasing or context sensitivity \\
Secure SDLC leadership & {[}B{]} 19--22; {[}D{]} 4; {[}C{]} 3 &
Assessment maturity and verification-gate evidence \\
\end{longtable}

These are electives. The baseline 16-week path deliberately does not
require reading all 3,000+ source pages cover to cover. Source chapters
deepen the specific competency, while the standalone guide and lab
evidence define measurable assessment outcomes.

\section{Annex C: deployment into a team learning
program}\label{annex-c-deployment-into-a-team-learning-program}

\textbf{Onboarding:} use Weeks 1--6 as an entry path for software
engineers entering AppSec. \textbf{Practitioner development:} pair Weeks
7--14 with real, explicitly authorized internal repository tasks;
protect customer data and keep incident response outside the training
sandbox. \textbf{Advanced engineers:} require deeper work on code-query
modeling, assurance cases, security policy composition, and architecture
review rather than simply increasing the number of alerts processed.

A small cohort benefits from a weekly 45-minute seminar, a 90-minute
supervised lab review, and asynchronous evidence submission. Assign
peers to defend or challenge source-to-sink traces, but grade based on
objective test evidence, correctness, and clear limits. Provide
alternate lab environments to avoid disadvantaging learners without
private-repo licenses.

\subsection{Change-control policy for the curriculum
itself}\label{change-control-policy-for-the-curriculum-itself}

\begin{enumerate}
\def\labelenumi{\arabic{enumi}.}
\tightlist
\item
  Keep the three source bibliographies stable and distinguish original
  text concepts from editorial interpretation.
\item
  Record tested tool versions, GitHub workflow actions, supported
  languages, and licensing assumptions in the course release notes.
\item
  Check OWASP ASVS/Top 10 and CWE pages for version changes before each
  cohort; do not reuse unversioned numeric requirements.
\item
  Review all sample code for language-version changes and obsolete
  framework recommendations.
\item
  Compile all eight LaTeX sources in a clean pdfLaTeX environment and
  inspect high-risk pages containing tables or listings.
\item
  Maintain an errata log: issue, affected module, cause, corrected
  source, and next validation date.
\end{enumerate}

\subsection{Suggested credential-independent graduation
record}\label{suggested-credential-independent-graduation-record}

\begin{Shaded}
\begin{Highlighting}[]
\NormalTok{Course version:       SSE{-}ASAS 1.0}
\NormalTok{Learner identifier:  [internal training ID]}
\NormalTok{Prerequisites:       verified}
\NormalTok{Lab IDs completed:  [L01 ... L14]}
\NormalTok{Mastery rubric:      [scores per competency]}
\NormalTok{Capstone SHA:        [immutable code reference]}
\NormalTok{Assessment reviewer: [person]}
\NormalTok{Outstanding gaps:    [explicit limitations]}
\NormalTok{Completion decision: [pass / revise]}
\end{Highlighting}
\end{Shaded}

This is an internal competency record, not a professional certification.
It makes progress auditable and prevents a calendar-based completion
claim from replacing demonstrated engineering competence.

\section{Reading links and source
register}\label{reading-links-and-source-register}

\textbf{{[}B{]}} Bishop, \emph{Computer Security: Art and Science}, 2nd
ed., provided PDF; focus Part II (Foundations), III (Policy), V
(Systems), VI (Assurance), VII (Special Topics).

\textbf{{[}D{]}} Dowd, McDonald and Schuh, \emph{The Art of Software
Security Assessment}, provided PDF; focus Part I (Introduction), Part II
(Software Vulnerabilities), Part III (Software Vulnerabilities in
Practice).

\textbf{{[}C{]}} Chess and West, \emph{Secure Programming with Static
Analysis}, provided PDF; focus Part I (Software Security and Static
Analysis), Part II (Pervasive Problems), Part III (Features and
Flavors), Chapters 13--14 exercises.

\textbf{{[}M1{]}} MITRE CWE \url{https://cwe.mitre.org/};
\textbf{{[}M2{]}} OWASP ASVS \url{https://owasp.org/projects/asvs};
\textbf{{[}M3{]}} OWASP Top 10:2025 \url{https://top10.owasp.org/2025/};
\textbf{{[}M4{]}} OWASP SAMM \url{https://owasp.org/projects/samm};
\textbf{{[}M5{]}} NIST SSDF
\url{https://csrc.nist.gov/pubs/sp/800/218/final}; \textbf{{[}M6{]}}
GitHub CodeQL CLI
\url{https://docs.github.com/en/code-security/concepts/code-scanning/codeql/codeql-cli};
\textbf{{[}M7{]}} GitHub Advanced Security product naming
\url{https://docs.github.com/en/get-started/learning-about-github/about-github-advanced-security}.

\end{document}
